\section{The least scalar profile}
\label{sec:profile}

We first isolate the scalar problem. The connection to actual tensors will
be proved by direct maps in Sections~\ref{sec:maps}--\ref{sec:extraction}.
In particular, no assumption that the least scalar profile is a tensor
character enters the construction.

\begin{definition}
A profile is a positive symmetric function \(P:\N_{>0}^2\to\mathbb R\)
satisfying
\begin{align}
 P(1,b)&=b,\label{eq:normalization}\\
 2P(a,b)&\ge P(a,b-1)+P(a,b+1)\quad(b\ge2),\label{eq:concavity}\\
 P(a,3h+a-1)&\ge3P(a,h)\quad(h\ge1).\label{eq:tripling}
\end{align}
\end{definition}

Define
\begin{equation}\label{eq:profile-formula}
 s_m=\prod_{h=1}^{m-1}\left(1+\frac1{3h}\right),\qquad
 F(m,n)=s_{\min(m,n)}\left(\max(m,n)+\frac{\min(m,n)-1}{2}\right).
\end{equation}
An empty product equals one.

\begin{theorem}[Least profile]\label{thm:least-profile}
The function \(F\) is a profile, and every profile \(P\) satisfies
\(P(m,n)\ge F(m,n)\). Moreover,
\begin{equation}\label{eq:profile-upper}
 F(m,n)\le(m+n-1)^{4/3},\qquad
 \log F(a,a)=\frac43\log a+O(1).
\end{equation}
The first inequality is strict when \(m=n>1\).
\end{theorem}
\begin{proof}
Normalization and symmetry are immediate. Using
\(s_{b+1}/s_b=(3b+1)/(3b)\) and
\(s_{b-1}/s_b=(3b-3)/(3b-2)\), direct subtraction gives
\[
 2F(a,b)-F(a,b-1)-F(a,b+1)=
 \begin{cases}
 \dfrac{2(a-b)s_b}{3b(3b-2)},&2\le b\le a,\\[4pt]
 0,&b>a.
 \end{cases}
\]
Thus each row is concave. Its affine tail has slope \(s_a\), so its
extension bounds the row from above:
\(F(a,h)\le s_a(h+(a-1)/2)\). Evaluating this affine tail at
\(3h+a-1\ge a\) proves \eqref{eq:tripling} for \(F\).

For minimality, write \(D_a=P(a,a)\),
\(k_a=2D_a/(3a-1)\), and
\(\Delta_{a,h}=P(a,h+1)-P(a,h)\).
Tripling at \(h=a\) and decreasing increments imply
\[
 2D_a\le P(a,4a-1)-P(a,a)\le(3a-1)\Delta_{a,a}.
\]
By symmetry and concavity, for \(a\ge2\),
\[
 D_a-D_{a-1}
 =\Delta_{a-1,a-1}+\Delta_{a,a-1}\ge k_{a-1}+k_a.
\]
It follows that \(k_a\ge(3a-2)k_{a-1}/(3a-3)\).
Since \(k_1=1\), this proves \(k_a\ge s_a\) and the diagonal bound.

Fix \(m<n\). Starting with \(x_0=m\), put
\(x_{j+1}=3x_j+m-1\), and choose \(j\) with \(x_j\ge n\).
Then
\[
 x_j-m=(3^j-1)(3m-1)/2,\qquad P(m,x_j)\ge3^jD_m.
\]
A concave sequence lies above the chord joining two of its entries.
The chord from \(m\) to \(x_j\) therefore gives
\[
 P(m,n)\ge D_m+\frac{n-m}{x_j-m}(3^j-1)D_m
 =\frac{2D_m}{3m-1}\left(n+\frac{m-1}{2}\right)\ge F(m,n).
\]
Only finitely many profile inequalities have been used.

Induction gives \(s_a^3\le(3a-1)/2\), strictly for \(a>1\).
The induction step reduces to
\((3a+2)(3a)^3-(3a-1)(3a+1)^3=6a+1>0\).
For \(m\le n\), both \((3m-1)/2\) and \(n+(m-1)/2\) are at most
\(m+n-1\). Cubing \eqref{eq:profile-formula} proves
\eqref{eq:profile-upper} and its strict diagonal version.
Finally, \(\log(1+x)=x+O(x^2)\) and convergence of
\(\sum h^{-2}\) give \(\log s_a=\frac13\log a+O(1)\).
\end{proof}

The theorem makes precise a limitation of the scalar argument: these
axioms, even together with \(P(a,b)\le(a+b-1)^{4/3}\), cannot force
diagonal growth faster than \(a^{4/3}\). This is a statement about the
scalar relaxation, not a claim that \(F\) comes from a tensor character.
The off-diagonal minimality uses shifted tripling instances with
\(h\) beyond the row index.
The diagonal bound and the exact certificate in Appendix~\ref{sec:exact-certificate} need only the
instances with \(h\) equal to the row index, which are the instances
realized by our tensor maps.


The affine rows of \(F\) will determine transition probabilities for the
finite construction. The scalar theorem identifies the best possible
diagonal growth under the stated axioms; no tensor character attaining
\(F\) is assumed.
